\documentclass[11pt]{book}
\usepackage{fontspec}
\setmainfont{Latin Modern Roman}
\setsansfont{Latin Modern Sans}
\usepackage{amsmath,amssymb,mathtools}
\usepackage{unicode-math}
\setmathfont{Latin Modern Math}
\usepackage{xurl}
\usepackage[hidelinks]{hyperref}
\newcommand{\sref}[2]{\hyperlink{source:#1:#2}{[S#2]}}
\newcommand{\cataloguescope}[1]{\paragraph{Mathematical question.}#1}
\begin{document}
\section{Research Problems in Function Theory — Problem 9.17}\label{statement}

\subsection{Definitions and mathematical statement}

\paragraph{Shared notation.}$\mathbb N=\{1,2,\ldots\}$, and $\mathbb Z,\mathbb Q,\mathbb R,\mathbb C$ denote the integers, rationals, reals and complex numbers. Any other conventions are specified locally.


Let $K\subset[0,1]$ be the middle-thirds Cantor set, equivalently \[K=\left\{\sum_{j=1}^{\infty}\frac{2\varepsilon_j}{3^j}:\varepsilon_j\in\{0,1\}\right\}.\] Identify $E=K\times K$ with $\{x+iy:x,y\in K\}\subset\mathbb C$, and let $\Omega=\widehat{\mathbb C}\setminus E$. The algebra $H^{\infty}(\Omega)$ consists of bounded analytic functions on this sphere domain, including analyticity at infinity.
\paragraph{Statement recorded in the supplied source.}
Prove the corona theorem for $\Omega$: for every finite $m\ge1$ and $f_1,\ldots,f_m\in H^{\infty}(\Omega)$ such that \[\inf_{z\in\Omega}\sum_{j=1}^m|f_j(z)|>0,\] there must exist $g_1,\ldots,g_m\in H^{\infty}(\Omega)$ with $\sum_{j=1}^mf_jg_j=1$ on $\Omega$. \sref{statement}{2}
\subsection{Short English statement}
Prove the bounded-analytic corona theorem for the complement on the sphere of the Cartesian square of the middle-thirds Cantor set.
\subsection{Sources}
\begin{description}
\item[\hypertarget{source:statement:1}{S1}] ULAM.AI and the UnsolvedMath Contributors, \emph{UnsolvedMath: A Curated Collection of Open Mathematics Problems}, catalogue entry. CC BY 4.0. \url{https://huggingface.co/datasets/ulamai/UnsolvedMath}.
\item[\hypertarget{source:statement:2}{S2}] Walter K. Hayman and Edward F. Lingham, \emph{Research Problems in Function Theory} (2018), arXiv:1809.07200. Problem 9.17, complete original question, all following updates, and the relevant chapter definitions. \url{https://arxiv.org/abs/1809.07200}.
\end{description}
\label{end:statement}
\end{document}
